% AI-written, AI cross-reviewed; no human sign-off; not a claim that RH is proved
% RH Notes blueprint source excerpts for Z1.1.06
% Verbatim excerpts in the notation of the parent blueprint. The labelled editions and source ranges appear below.
% See the unit YAML for the current proof status. Earlier-round passages retain their original wording.
\section*{Z1.1.06. The common marking and degree zero}

% Parent source SHA-256: 9559fca8cc392b75d71f1c158fe0c635cfcab567e2f489b52db4c5bb29d98029
% blueprint/ch_P_r8_small.tex lines 458 through 479; r8:small:Z1.1.06
\paragraph{Z1.1.06. Detecting the ramification of the quadratic character.}
\phantomsection\label{r8:small:Z1.1.06}
Let \(A\) be a valuation ring in the chosen algebraic closure and \(K\) the
quadratic field whose stabilizer is \(\ker\eta\) in the marking data.
Each image of \(\mathcal O_K\) is contained in \(A\) by S10.
Suppose \(\mathrm{disc}(K)\) has nonzero residue. In any integral basis
of \(\mathcal O_K\), the square of the determinant of the matrix of its
embeddings is that discriminant. If two distinct embeddings agreed modulo
the maximal ideal on the basis, reduction would give equal rows, a zero
determinant and a zero discriminant. Thus inertia, which fixes each
residue, fixes the restriction to \(K\) of the given embedding. In
particular, if inertia is not contained in \(\ker\eta\), the discriminant
is in the maximal ideal of \(A\). Its absolute integer value is there
as well, and multiplying by the integral \(\ell N\) keeps it there.
This proves the ramification test at the enlarged level. S06 explains
inversion of the algebraic twist on the same Artin map. The linked
endpoint data and descent giving the two markings remain to be checked
(\nolinkurl{GlobalTheory.UnramDiscr.embedding_eq},
\path{IneLeFixingSubgroup.lean:18--58};
\nolinkurl{gdEnd_of_frame_linkedEnd_downCsdK_descent},
\path{DenseThresholdCases/GdEndOf.lean:54--83};
\nolinkurl{ends_of_alt}, \path{DenseThresholdCases/DnEnds.lean:35--39}).

